\documentclass[10pt,a4paper]{amsart}
\usepackage[margin=19mm]{geometry}
\usepackage{amsmath,amssymb,mathtools}
\usepackage[scr=rsfs,cal=euler]{mathalfa}
\usepackage{xfrac}
\usepackage[unicode,hidelinks,bookmarks=false]{hyperref}
\newcommand{\ie}{i.e.\ }
\newcommand{\cf}{cf.\ }
\newcommand{\resp}{respectively\ }
\newcommand{\Subsection}{\subsection}
\providecommand{\nobreakdash}{\nobreak\hbox{-}}
\setlength{\emergencystretch}{2em}
\multlinegap0pt
\usepackage{bm}
\usepackage{bbm}
\usepackage{dsfont}
\usepackage{xspace}
\usepackage{cancel}
\usepackage{mathtools}
\usepackage{amsthm}
\makeatletter
\@ifundefined{theorem}{\newtheorem{theorem}{Theorem}}{}
\@ifundefined{proposition}{\newtheorem{proposition}[theorem]{Proposition}}{}
\makeatother
\title[Hidden convexity of quadratic systems and its application to]{Hidden convexity of quadratic systems and its application to quadratic programming}
\author{Nguyen Quang Huy, Nguyen Huy Hung, Tran Van Nghi, Hoang Ngoc Tuan, Nguyen Van Tuyen}
\date{Source: arXiv:2601.13511v1 (CC BY 4.0)}
\begin{document}
\maketitle

\noindent\emph{Source and licence.} Hidden convexity of quadratic systems and its application to quadratic programming. Version arXiv:2601.13511v1 (CC BY 4.0), distributed under the Creative Commons Attribution 4.0 International licence, as recorded in the arXiv licence metadata for this exact version and shown on the abstract page https://arxiv.org/abs/2601.13511v1. Full rights notice: licenses/arxiv-2601.13511-CC-BY-4.0.txt.\par\smallskip

\noindent\textit{Editorial scope.} Counted unit: the Proposition whose source label is \texttt{pro1} in the article (source0.tex lines 1084--1099, source label \texttt{pro1}), delivered together with the article's own complete original proof of it (source0.tex lines 1100--1139, one \texttt{proof} environment of 40 lines). No mathematics is written, repaired or re-derived: statement and proof are the article's own text line for line. Required context retained uncounted: Theo7 (lines 806--842). Every definition, display and label of those lines is retained unchanged. Omitted from this package and not redistributed: 1--805, 843--1083, 1140--1323. Nothing inside a retained excerpt is omitted or altered: every retained body is delivered exactly as the article's lines are written, so no omission receipt of any kind accompanies this package. Citations: the retained excerpts carry no citation command at all, so no bibliography entry of the article is transcribed and no reference is redistributed. Reference adaptation: none was needed and none was made; every reference inside the delivered unit resolves inside it, so no unresolved reference or citation is produced. Printed slips kept literal: no printed word, symbol, hypothesis or number of the counted statement or of its proof was altered. Layout and class: the article's own class is \texttt{interact} with packages including epstopdf, subfigure, hyperref, natbib; the wrapper is the lane's pinned \texttt{amsart} editorial wrapper at 10pt on A4, which restates the theorem-like environments the retained text needs and adds the article's own macro definitions that the retained text needs (none), with extra wrapper packages (bm, bbm, dsfont, xspace, cancel, mathtools). The delivered numbering is the unit's own. Assets: no retained span contains a figure environment, a graphics-inclusion command, a PDF-page inclusion, a table or a TikZ picture; no image file is copied into this package, the package declares no asset dependency and no image file is redistributed with it. Licence: version arXiv:2601.13511v1 of this article is distributed under the Creative Commons Attribution 4.0 International licence, as recorded in the arXiv licence metadata for this exact version and shown on the abstract page https://arxiv.org/abs/2601.13511v1. A full rights notice is delivered at licenses/arxiv-2601.13511-CC-BY-4.0.txt. Characters: every retained excerpt is pure ASCII, so the delivered engine, XeLaTeX with its default Latin Modern text font, reports no missing character.
	\begin{theorem}\label{Theo7} 
		Consider the following statements:
		\begin{description}
			\item{$(i)$} {\rm [S-lemma]} For each quadratic function $f:\mathbb{R}^n\to \mathbb{R}$,
			\begin{align*}
				&[g_i(x)\leq 0,\; i=1, \ldots, m \Rightarrow f(x)\geq 0] \\[0.3em]
				\Leftrightarrow \; \; &(\exists \lambda_i\geq0,\; i=1, \ldots, m)\; (\forall x\in \mathbb{R}^n)\;  
				f(x)+\sum_{i=1}^{m}\lambda_ig_i(x)\geq 0;
			\end{align*}
			
			\item{$(ii)$} The set
			$$\bigcup_{\lambda\in  \mathbb{R}^m_+} \; {\rm epi}\bigg(\sum_{i=1}^{m}\lambda_ig_i\bigg)^*$$
			is closed;
			
			\item{$(iii)$} {\rm [Slater's condition]} There exists $x^0\in \mathbb{R}^n$ such that
			$g_i(x^0)< 0,\; i=1, \ldots, m$;
			
			\item{$(iv)$} {\rm [Strong duality]} For each quadratic function $f$,
			\begin{equation}\label{e11} 
				\inf \{f(x): g_i(x)\leq 0,\; i=1,\ldots,m\}
				=\max_{\lambda\in \mathbb{R}^m_+}\inf_{x\in \mathbb{R}^n}
				\bigg\{f(x)+\sum_{i=1}^{m}\lambda_ig_i(x)\bigg\}.
			\end{equation}
		\end{description}
		The following statements hold:
		\begin{description}
			\item{$(a)$} $(i)$ implies $(ii)$;
			
			\item{$(b)$} If ${\rm U}(f, g_1, g_2,\ldots,g_m)$ is a convex set, then $(iii)$ implies $(i)$; 
			
			\item{$(c)$} If the assumptions $({\rm H_1})$ and $({\rm H_2})$ hold, then $(iii)$ implies $(i)$;
			
			\item{$(d)$} $(i)$ and $(iv)$ are equivalent;
			
			\item{$(e)$} If $m=1$, then the statements $(i)$--$(iv)$ are equivalent.
		\end{description}
	\end{theorem}

	\begin{proposition}\label{pro1}
		Consider the following quadratic programming problem:
		\begin{equation}\tag{$QP$}
			\inf\{f(x): g_i(x)\leq 0,\; i=1, \ldots, m\}.
		\end{equation}
		Suppose that assumptions $({\rm H_1})$ and $({\rm H_2})$ are satisfied, and that there exists
		$x^0\in \mathbb{R}^n$ such that $g_i(x^0)<0$ for every $i=1, \ldots, m$.
		Let $\bar x$ be a feasible point of problem $(QP)$.
		Then, $\bar x$ is a global minimizer of $(QP)$ if and only if there exists
		$\lambda=(\lambda_1,\ldots,\lambda_m)\in \mathbb{R}^m_+$ such that the following conditions hold:
		\begin{description}
			\item{$(i)$} $2(A+\sum_{i=1}^{m}\lambda_iB_i)\bar x+ (a+\sum_{i=1}^{m}\lambda_ib_i)=0$;
			\item{$(ii)$} $\lambda_ig_i(\bar x)=0, \; i=1, \ldots,m$; and,
			\item{$(iii)$} $A+\sum_{i=1}^{m}\lambda_iB_i\succeq 0.$
		\end{description}
	\end{proposition}

	\begin{proof}
		According to Theorem~\ref{Theo7}, $\bar x$ is a global minimizer of $(QP)$ if and only if there exists
		$\lambda=(\lambda_1,\ldots,\lambda_m)\in \mathbb{R}^m_+$ such that
		\begin{equation}\label{e12}
			f(x)+\sum_{i=1}^{m}\lambda_i g_i(x)\geq f(\bar x), \qquad \forall x\in \mathbb{R}^n.
		\end{equation}
		From this inequality and the fact that $\bar x$ is a feasible point of $(QP)$, it follows that
		\begin{equation}\label{e13}
			\sum_{i=1}^{m}\lambda_i g_i(\bar x)=0.
		\end{equation}
		Since $\lambda_i\geq 0$ and $g_i(\bar x)\leq 0$ for every $i=1,\ldots,m$, we obtain
		\begin{equation}\label{e14}
			\lambda_i g_i(\bar x)=0, \qquad i=1,\ldots,m.
		\end{equation}
		
		Define
		$$\varphi(x):=f(x)+\sum_{i=1}^{m}\lambda_i g_i(x)-f(\bar x).$$
		Since $\bar x$ is a global minimizer of the unconstrained problem
		$\inf\{\varphi(x): x\in \mathbb{R}^n\}$, it follows that
		$$\nabla \varphi(\bar x)=0
		\quad \text{and} \quad
		\nabla^2 \varphi(\bar x)\succeq 0,$$
		that is,
		$$2\bigg(A+\sum_{i=1}^{m}\lambda_i B_i\bigg)\bar x
		+\bigg(a+\sum_{i=1}^{m}\lambda_i b_i\bigg)=0$$
		and
		$$A+\sum_{i=1}^{m}\lambda_i B_i\succeq 0.$$
		
		Conversely, suppose that conditions $(i)$--$(iii)$ are satisfied.
		Then $\bar x$ is a global minimizer of the convex quadratic programming problem
		$\inf\{\varphi(x): x\in \mathbb{R}^n\}$.
		Consequently, for all $x\in \mathbb{R}^n$, we have
		\begin{align*}
			f(x)
			&\geq f(x)+\sum_{i=1}^{m}\lambda_i g_i(x) \\
			&\geq f(\bar x)+\sum_{i=1}^{m}\lambda_i g_i(\bar x)
			= f(\bar x),
		\end{align*}
		which shows that $\bar x$ is a global minimizer of $(QP)$.
	\end{proof}

\end{document}
